%%
%% Chapter 1: Introduction
%%
\chapter{Introduction}
\label{ch:01}

\section{Motivation}
\label{sec:intro-motivation}

Video surveillance has become one of the most widely deployed sensing technologies in public and
private spaces. Transportation hubs, retail environments, campuses, and city streets are now covered
by camera networks whose combined output far exceeds what human operators can monitor: a single
operator can attend to only a handful of feeds at a time, and attention degrades rapidly over
sustained viewing. As the scale of these systems continues to expand, automated analysis of
surveillance footage has shifted from a convenience to a necessity~\cite{sultani2018ucfcrime}.

Within the space of events that automated systems should detect, violence-related anomalies occupy a
special position. Fights, assaults, robberies, shootings, and road accidents are precisely the
incidents for which a delayed response carries the highest human cost, and they are also
comparatively rare---a property that makes them poor targets for conventional supervised
classification, which requires abundant labeled examples of every class. Video anomaly detection
(VAD) reframes the problem: rather than recognizing a fixed vocabulary of actions, the system
assigns each frame an \emph{anomaly score} indicating how strongly it deviates from normal activity,
so that rare and heterogeneous dangerous events can be surfaced without exhaustive per-class
supervision.

The central practical obstacle for VAD is annotation cost. Training a frame-level detector directly
would require annotators to mark the exact temporal extent of every anomalous event in every
training video---a task that is expensive, subjective at event boundaries, and infeasible at the
scale of modern datasets containing thousands of hours of footage. The \emph{weakly supervised}
formulation addresses this by requiring only \emph{video-level} labels: each training video is
tagged as containing an anomaly somewhere, or as entirely normal. This paradigm, established by
Sultani et al.~\cite{sultani2018ucfcrime} together with the UCF-Crime benchmark, casts training as
multiple instance learning (MIL): a video is a \emph{bag} of temporal segments (instances), the bag
label is known, but the instance labels are not. A ranking loss encourages the highest-scoring
segments of an anomalous bag to score above the highest-scoring segments of a normal bag, allowing
the network to localize anomalies at test time despite never seeing frame-level supervision during
training. Subsequent refinements such as Robust Temporal Feature Magnitude learning
(RTFM)~\cite{tian2021rtfm} strengthened this formulation, and weakly supervised VAD (WSVAD) has
since become the dominant paradigm on the two principal violence-oriented benchmarks,
UCF-Crime~\cite{sultani2018ucfcrime} and XD-Violence~\cite{wu2020xdviolence}.

A second transformative development is the arrival of vision-language models (VLMs). Contrastively
trained image encoders such as CLIP~\cite{radford2021clip} and the sigmoid-loss SigLIP
family~\cite{zhai2023siglip,tschannen2025siglip2} produce semantically rich frame representations
that transfer remarkably well to downstream tasks without fine-tuning. Methods such as
VadCLIP~\cite{wu2024vadclip} demonstrated that \emph{frozen} CLIP features, combined with
lightweight temporal modeling, are highly effective for WSVAD. Frozen-feature designs are especially
attractive in resource-constrained research and deployment settings: the expensive backbone is run
once per video as a feature extractor, and only a small head is trained.

Yet appearance is not the whole story. Violence is fundamentally a phenomenon of \emph{motion}---the
rapid, forceful, articulated body dynamics of fighting, striking, or falling. Skeleton-based action
recognition, exemplified by graph convolutional networks over human pose sequences such as
CTR-GCN~\cite{chen2021ctrgcn}, captures exactly this signal, is invariant to many appearance
nuisances (clothing, background, lighting), and has driven a distinct line of skeleton-only
anomaly detection---from reconstructing normal pose
trajectories~\cite{morais2019skeleton}, through pose-graph
clustering~\cite{markovitz2020gepc}, to normalizing-flow density
models~\cite{hirschorn2023stgnf}. This complementarity---semantic appearance from a
VLM, articulated dynamics from a pose graph---motivates the dual-modal design studied in this
thesis.

Finally, deployment realism demands attention to \emph{robustness}. A model trained on curated
benchmark footage will, in the field, encounter sensor noise from low-light cameras, motion blur,
aggressive compression from bandwidth-limited transmission, and illumination changes. Test-time
adaptation (TTA) methods such as TENT~\cite{wang2021tent} and SAR~\cite{niu2023sar} promise to
recover accuracy under such distribution shift without labels, by adapting normalization-layer
parameters through entropy minimization. Whether these promises survive contact with a frozen,
LayerNorm-based~\cite{ba2016layernorm}, rank-evaluated VAD architecture is one of the questions this
thesis answers---in the negative---before proposing an alternative that does work.

\section{Problem Statement}
\label{sec:intro-problem}

This thesis addresses weakly supervised violence-oriented video anomaly detection under the
following setting. Given a training corpus of untrimmed videos, each carrying only a binary
video-level label (anomalous somewhere / entirely normal), the goal is to learn a scoring function
that, at test time, assigns each frame of an unseen video an anomaly score. Performance is measured
by frame-level ranking metrics against held-out temporal annotations: the area under the ROC curve
(AUC) on UCF-Crime and average precision (AP) on XD-Violence, following each benchmark's standard
protocol~\cite{sultani2018ucfcrime,wu2020xdviolence}.

Within this setting, the thesis studies three interlocking problems:

\begin{enumerate}
    \item \textbf{Dual-modal fusion.} How should a frozen skeleton encoder (CTR-GCN over
RTMPose~\cite{jiang2023rtmpose} keypoints) and a frozen vision-language encoder be combined so that
the weaker-standalone skeleton stream \emph{adds} information rather than diluting the stronger
visual stream? The design space ranges from fixed score averaging (late fusion) to learned,
input-dependent gating.
    \item \textbf{Backbone sensitivity.} Frozen-feature pipelines inherit the properties of their
backbone. With four candidate visual backbones spanning two architectural families and an order of
magnitude in capacity---CLIP ViT-B/16, SigLIP2 ViT-B/16, SigLIP2 SO400M, and SigLIP2 Giant---does
the ranking of backbones transfer across datasets, and does fusion interact with backbone choice?
    \item \textbf{Test-time robustness.} Under synthetic but realistic input corruptions, can the
frozen dual-modal detector be adapted at test time---without labels, and without tuning
hyperparameters on test data---to recover the accuracy lost to corruption? In particular, do
entropy-minimization TTA methods designed for BatchNorm~\cite{ioffe2015batchnorm} networks transfer
to a LayerNorm-based fusion head evaluated by a rank metric?
\end{enumerate}

All experiments are constrained to a single NVIDIA RTX 4090 GPU (24~GB), which rules out end-to-end
backbone fine-tuning and motivates the frozen-feature design throughout: only a lightweight fusion
head (0.50M--1.02M trainable parameters, depending on backbone) is learned, and each configuration
trains in under five minutes. This constraint is not merely practical; it defines the regime the
thesis characterizes---what a reproducible, single-GPU, frozen-feature dual-modal system can and
cannot achieve.

\section{Research Gaps}
\label{sec:intro-gaps}

Four gaps in the literature motivate this work.

\textbf{Gap 1: Single-modality dominance in WSVAD.} The overwhelming majority of weakly supervised
VAD methods score a single visual feature stream---C3D or I3D motion features in the earlier
generation, frozen CLIP features in the current one. Skeleton-based anomaly detection exists as a
separate, largely unsupervised line of
work~\cite{morais2019skeleton,markovitz2020gepc,hirschorn2023stgnf}, but the integration of
skeleton dynamics with vision-language features in a unified weakly supervised framework remains
largely unexplored.
Whether pose dynamics carry signal that frozen VLM features miss---and whether that signal survives
fusion with a much stronger visual stream---is an open empirical question.

\textbf{Gap 2: CLIP-era methods overlook motion dynamics.} Vision-language backbones are trained on
static image--text pairs; their per-frame features encode scene semantics, objects, and context, but
only weakly encode articulated human motion. Violence-related anomalies---fighting, assault,
vehicular collisions---are characterized by distinctive spatiotemporal motion patterns that static
appearance features may not fully capture. CLIP-based WSVAD methods inherit this blind spot, and
none of them measure what an explicit motion modality would add.

\textbf{Gap 3: Backbone sensitivity is unstudied.} Nearly every published WSVAD method reports
results with one fixed backbone, implicitly assuming its conclusions are backbone-independent. The
rapid proliferation of VLM variants---CLIP, SigLIP, SigLIP2 at multiple
scales~\cite{radford2021clip,zhai2023siglip,tschannen2025siglip2}---makes this assumption untenable:
capacity, training objective, and training data all differ. No prior work systematically holds the
WSVAD pipeline fixed and varies the visual backbone across two benchmarks, let alone examines how
backbone choice interacts with multi-modal fusion.

\textbf{Gap 4: Entropy-based TTA assumes BatchNorm; modern heads use LayerNorm.}
TENT~\cite{wang2021tent} and SAR~\cite{niu2023sar} were designed for networks whose BatchNorm layers
maintain running statistics of the source distribution; adaptation recalibrates those statistics and
their affine transforms. Transformer-era architectures---including VLM backbones and the fusion head
studied here---use LayerNorm~\cite{ba2016layernorm}, which normalizes per sample and keeps no
running statistics. Whether entropy minimization confined to LayerNorm affine parameters can help a
frozen, rank-evaluated detector has, to our knowledge, not been investigated in the VAD literature;
nor has the question of what \emph{does} restore robustness when it cannot.

\section{Research Questions}
\label{sec:intro-rqs}

The gaps above condense into three research questions.

\begin{description}
    \item[RQ1 (Dual-modal fusion value).] Does fusing frozen CTR-GCN skeleton features with frozen
vision-language features improve weakly supervised violence detection over the visual-only
baseline---and which fusion mechanism (fixed late fusion vs.\ learned gated fusion) is required for
the skeleton stream to help rather than hurt?
    \item[RQ2 (Backbone sensitivity).] How sensitive are fusion performance and its conclusions to
the choice of visual backbone across CLIP ViT-B/16, SigLIP2 ViT-B/16, SigLIP2 SO400M, and SigLIP2
Giant---and are backbone preferences consistent between UCF-Crime and XD-Violence?
    \item[RQ3 (Entropy-TTA transfer to LayerNorm).] Do entropy-minimization TTA methods (TENT, SAR)
transfer from their native BatchNorm setting to the LayerNorm-based fusion head of a frozen
rank-evaluated VAD model under input corruption---and if not, what label-free, tuning-free
alternative recovers corrupted-input accuracy?
\end{description}

Chapters~\ref{ch:05} and~\ref{ch:06} answer these questions empirically: RQ1 and RQ2 in
Chapter~\ref{ch:05} through a six-variant, four-backbone, three-seed ablation study on both
benchmarks, and RQ3 in Chapter~\ref{ch:06} through a 20-condition corruption study with a full TTA
method comparison.

\section{Contributions}
\label{sec:intro-contributions}

This thesis makes three contributions. All results rest on a uniform three-seed
experimental protocol spanning six fusion variants and four backbones
(Tables~\ref{tab:ucf-ablation} and~\ref{tab:xd-ablation}), complemented by
per-category and failure analysis (Sections~\ref{sec:qualitative}--\ref{sec:percat})
and a complete test-time-adaptation investigation---the six-strategy search, the
NORM/CORAL eliminations, and the episodic-inertness diagnosis
(Chapter~\ref{ch:06}, Section~\ref{sec:coral}).

\textbf{1. A dual-modal skeleton--visual fusion framework for WSVAD.} We propose a pipeline that
combines frozen CTR-GCN skeleton features with frozen vision-language features through learned gated
fusion, trained with MIL ranking loss. The gated mechanism projects both modalities to a shared
space, computes a per-dimension gate from their concatenation, and blends them over a residual
connection that always retains both streams---so the gate modulates each modality's \emph{relative}
contribution rather than suppressing either outright. Only the fusion head and scoring head are
trained; both backbone encoders remain frozen. The framework achieves 82.5$\pm$0.4\% AUC on
UCF-Crime (with SigLIP2 Giant) and 78.7$\pm$0.9\% AP on XD-Violence (with SigLIP2 SO400M), averaged
over three seeds. The skeleton stream's contribution is deliberately characterized honestly: gated
fusion meets or exceeds the visual-only baseline in all eight backbone$\times$dataset configurations
(mean $+0.6$ points; strictly positive in six of the eight, the other two tied at the displayed
precision)---a small but consistent complementary gain, not a transformative one. Equally
informative is the negative half of the finding: naive equal-weight late fusion frequently
\emph{degrades} performance below the visual-only baseline, indicating that a learned
combination---here, input-dependent gating---is what allows a weak modality to help.

\textbf{2. A systematic four-backbone study.} We evaluate the identical pipeline across four visual
backbones---CLIP ViT-B/16, SigLIP2 ViT-B/16, SigLIP2 SO400M, and SigLIP2 Giant---on both UCF-Crime
and XD-Violence, with three seeds per configuration. The study reveals dataset-specific backbone
preferences that single-backbone evaluations cannot see: UCF-Crime's fixed-camera surveillance
footage favors the largest backbone (Giant), whereas on XD-Violence's heterogeneous sources the four
backbones are statistically indistinguishable visual-only and SO400M shows the highest observed mean
only under gated fusion. The skeleton stream's complementary value itself varies with the backbone
it is fused with---a dependence invisible to single-backbone studies and directly relevant to
practitioners choosing a backbone from general-purpose benchmark rankings.

\textbf{3. A test-time adaptation study for frozen dual-modal VAD, with a working alternative.} On a
20-condition corruption benchmark (UCF-Crime-C: four corruption types at five severities), we show
that entropy-minimization TTA (TENT, SAR) applied to the fusion head's LayerNorm affine
parameters---1{,}536 scalars, the only adaptable surface in an otherwise frozen model---changes
corrupted AUC by less than $0.1$ points on every backbone. We argue this is structural rather than a
tuning artifact: the corruption shift lives upstream in the frozen backbone features, and affine
nudges of this size cannot reorder the pooled snippet rankings on which AUC depends. We further show
that naive feature-statistic restoration (NORM/CORAL-style) is rank-disruptive with a
backbone-dependent sign. We then propose \emph{discriminative-reliability reweighting}: a
label-free, tuning-free, zero-parameter-update method that detects when the visual-language stream's
discriminative signal collapses---via the dispersion of its solo anomaly scores---and routes the
gated fusion toward the corruption-robust skeleton stream, guarded by a skeleton-reliability gate.
This improves corrupted AUC on all four backbones (mean $+1.2$ points over three seeds, up to
$+13.2$ points in the condition that most collapses the visual-language stream), while leaving the
model bit-identical to the source model wherever adaptation would be unsafe.

\section{Scope and Honest Positioning}
\label{sec:intro-scope}

The results in this thesis should be read against an explicit statement of what they do and do not claim.

\textbf{No peak-score claim.} The headline numbers---82.5$\pm$0.4\% UCF-Crime AUC and 78.7$\pm$0.9\%
XD-Violence AP---trail the field's fine-tuned and text-aligned leaders, which reach roughly 88--91\%
UCF-Crime AUC. That gap follows from deliberate design choices: both backbones frozen, no
text-prompt branch, no temporal transformer, no auxiliary modalities, single-GPU training. The
contribution is a reproducible, single-GPU dual-modal pipeline with a complete ablation and
robustness analysis---not a new peak score. Chapter~\ref{ch:07} positions the system against the
field in detail, including a like-for-like comparison against methods operating under comparable
frozen-feature constraints.

\textbf{Robustness gains are corrupted-AUC gains.} The test-time reweighting method of
Chapter~\ref{ch:06} improves AUC \emph{under corruption} only; it leaves clean-data accuracy
untouched by construction. Its reliability statistics are pooled over each corruption condition (a
transductive assumption, standard in this benchmark setting but disclosed throughout); a streaming
variant is future work (Section~\ref{sec:conc-future}).

\textbf{Reproducibility as a first-class goal.} Every reported number derives from fixed seeds (42,
123, 2024), fixed train/validation splits, and a fixed evaluation protocol; results are reported as
mean$\pm$std over the three seeds. The complete configuration provenance---committed configs and
data manifests, and tracked per-run metrics files---backs every reported number.

\section{Thesis Organization}
\label{sec:intro-roadmap}

The remainder of this thesis is organized as follows. Chapter~\ref{ch:02} surveys related work
across the four lineages this thesis draws on: weakly supervised VAD from MIL ranking through the
CLIP era to current multi-modal systems, skeleton-based action recognition, vision-language
backbones, and test-time adaptation. Chapter~\ref{ch:03} presents the methodology: the dual-modal
architecture, both feature extraction pipelines, the fusion mechanisms, MIL training, and the
test-time adaptation methods including the proposed discriminative-reliability reweighting.
Chapter~\ref{ch:04} details the experimental setup---datasets, evaluation protocol, the UCF-Crime-C
corruption benchmark, implementation details, and infrastructure. Chapter~\ref{ch:05} reports the
fusion and backbone results (RQ1, RQ2), including ablations, per-category analysis, statistical
robustness checks, and failure analysis. Chapter~\ref{ch:06} reports the robustness and TTA results
(RQ3), from the entropy-minimization null result through the reweighting method's evaluation.
Chapter~\ref{ch:07} discusses the findings, positions the system against the state of the field,
and states the study's limitations; Chapter~\ref{ch:09} concludes and outlines future work.
